\section{Introduction}
Audio temporal grounding maps a query to event intervals~\cite{xu2021grounding,munakata2025amr}. In a recording containing both ``dog then rooster'' and ``rooster then dog,'' the descriptions refer to different windows. Recognizing the categories establishes presence; answering the query requires selecting the requested order. Returning the same real window for both queries is a \emph{query--window binding error}. We target this distinction between event detection and query-specific localization.

Grounding draws on weak-caption learning~\cite{xu2024wstag}, audio-language interaction in Pengi~\cite{deshmukh2023pengi}, Qwen2-Audio~\cite{chu2024qwen2audio}, and Audio Flamingo~3~\cite{goel2025af3}, temporal representations in TimeAudio~\cite{wang2026timeaudio}, and interleaved timestamps with absent-event training in SpotSound~\cite{sun2026spotsound}. We focus on competition between \emph{present} events: a plausible answer can contain the correct sounds in the wrong relation. This requires explicit comparison of competing local answers.

CLAP learns audio--text matching~\cite{wu2023clap}; T-CLAP combines concatenated sounds, reversed descriptions, and temporal contrastive training~\cite{yuan2024tclap}. CompA examines order and binding~\cite{ghosh2024compa}, CoSTALA learns hierarchical spatio-temporal alignment~\cite{ren2026costala}, and SHINE uses compositional negative queries for video grounding~\cite{cheng2024shine}. RelTwin instead compares local timestamp answers: both inverse descriptions are true in one recording, and the negative window answers the opposite query.

RelTwin (Fig.~\ref{fig:pair}) addresses this binding problem through \emph{locally valid timestamp contrasts}. Shared excerpts form opposite-order windows in one recording. Each window contains a real event sequence, but its correctness as an answer depends on the query. Sequence supervision rewards correct timestamp answers individually; the added candidate loss explicitly compares the two answers under each query. This local competition trains query--window binding while inference retains direct timestamp generation.

We make three contributions. First, we characterize query--window binding errors when inverse relations co-occur in one recording. Second, we introduce locally valid timestamp contrasts to supervise the query-specific choice between competing answers. Third, paired diagnosis and a timing intervention demonstrate improved query--window binding, alongside gains on SpotSound-Bench.
